For the Wick-ordered $\lambda\varphi^4$ measure in two dimensions, uniqueness of the invariant measure for the dynamics is known above the susceptibility threshold $\mu_c$. Whether the magnetisation is spontaneously non-zero for every $\mu<\mu_c$, the other half of sharpness, is open in the continuum. We do not settle it. We prove that for every even single-site measure $e^{-V}$ with $V'$ convex on $[0,\infty)$, the ratio $\langle\varphi;\varphi^2\rangle_H/\langle\varphi\rangle_H$ is nonincreasing in the field $H>0$, double wells included. On the lattice torus we prove the exact identity $m-h\chi=4\lambda m^2\partial_hR+8\lambda m^2(-\partial_\mu m)$. With a localised characteristics lemma of Aizenman--Barsky--Fern\'andez type, it reduces continuum sharpness, with an explicit lower bound on the spontaneous magnetisation, to one structural inequality (P). Inequality (P) is needed only at small magnetisation and on the field scale $h\lesssim1/\chi$, uniformly in the lattice spacing, and with any Dini modulus; the Dini condition is sharp. We give three exact reformulations of (P). We also record obstructions, among them that (P) fails for the two-dimensional plane rotator, so any proof must use one-component structure. Monte Carlo on two-dimensional tori in the phase-flipping regime finds the strongest form of (P) satisfied by a margin of at least $13\%$. The one-site theorem, the torus identities and the characteristics lemma are formalised in Lean~4 with Mathlib. The work was carried out by a multi-model swarm of AI agents on a public research workspace, under the direction of the author, who takes responsibility for its content.
